\documentclass[a4paper, 11pt]{article}

% --- Packages ---
\usepackage[utf8]{inputenc}
\usepackage[T1]{fontenc}
% \usepackage[french]{babel}
\usepackage[margin=1.5cm, top=2cm, bottom=2cm]{geometry}
\usepackage{xcolor}
\usepackage{fancyhdr}
\usepackage{titlesec}
\usepackage{enumitem}
\usepackage{array}
\usepackage{longtable}
\usepackage{booktabs}
\usepackage{colortbl}

% --- Définition des couleurs ---
\definecolor{primaryBlue}{RGB}{20, 50, 90}
\definecolor{lightBlue}{RGB}{235, 240, 247}

% --- Styles de document ---
\pagestyle{fancy}
\fancyhf{}
\fancyfoot[L]{\footnotesize \textcolor{gray}{Projet Wakala - One Pager (Résumé)}}
\fancyfoot[R]{\footnotesize \textcolor{gray}{Page \thepage}}
\renewcommand{\headrulewidth}{0pt}
\renewcommand{\footrulewidth}{0.4pt}

% --- Personnalisation des titres ---
\titleformat{\section}{\Large\bfseries\color{primaryBlue}}{}{0em}{}[\titlerule]
\titleformat{\subsection}{\large\bfseries\color{primaryBlue}}{}{0em}{}

\newcolumntype{L}[1]{>{\raggedright\arraybackslash}p{#1}}

\begin{document}

% --- En-tête ---
\begin{center}
    {\Huge \textcolor{primaryBlue}{\textbf{WAKALA}}}\\[0.2cm]
    {\large \bfseries \textcolor{primaryBlue}{Résumé du Projet (One Pager)}}
\end{center}

\vspace{0.5cm}

% --- Introduction ---
\section*{Introduction}
L'objectif de ce projet de stage est de concevoir un \textbf{Chatbot interactif couplé à un Moteur de Recommandation intelligent} pour guider et aider chaque client à trouver la voiture parfaite pour lui, de A à Z. Wakala ne se contente plus de lister des voitures, elle devient un véritable conseiller virtuel personnalisé.

\vspace{0.3cm}

% --- Contenu ---
\section*{1. Le Challenge et la Stratégie}
Le défi principal de l'achat automobile en ligne au Maroc est le manque d'accompagnement et la rigidité des filtres classiques. La stratégie de Wakala pour y répondre s'articule autour de l'Intelligence Artificielle pour remplacer la recherche par un véritable accompagnement personnalisé.

\section*{2. Périmètre et Fonctionnalités Attendues}
Voici notre plan d'action pour créer cette plateforme d'accompagnement intelligent :
\begin{itemize}[label=\textbullet]
    \item \textbf{Comprendre avant de proposer :} Un agent conversationnel (Chatbot) qui discute naturellement avec l'utilisateur pour définir un profil acheteur précis (budget, usage, famille).
    \item \textbf{Recommander intelligemment :} Un moteur d'algorithmes croise ce profil avec le marché pour suggérer les véhicules parfaits, au lieu d'afficher une simple liste.
    \item \textbf{Être totalement transparent :} L'IA explique exactement \textit{pourquoi} une voiture est recommandée, pour instaurer la confiance avec l'acheteur.
    \item \textbf{Guider l'acheteur pas à pas :} L'assistant virtuel accompagne l'utilisateur de la définition du besoin jusqu'au choix final du véhicule, de manière fluide et rassurante.
\end{itemize}

\section*{3. Ressources et Moyens Alloués}
Pour réussir ce projet de manière agile et innovante, nous nous appuyons sur les ressources suivantes :

\vspace{0.2cm}
\textbf{Ressources Humaines}
\begin{itemize}[label=\textbullet]
    \item \textbf{Une équipe technique spécialisée :} Des développeurs et experts en intelligence artificielle (Machine Learning, NLP) pour concevoir les algorithmes et le chatbot.
    \item \textbf{Une approche centrée utilisateur :} Des concepteurs d'interface (UX/UI) pour s'assurer que la conversation et les recommandations sont claires et faciles à utiliser.
\end{itemize}

\textbf{Ressources Technologiques}
\begin{itemize}[label=\textbullet]
    \item \textbf{Des modèles génératifs puissants :} L'utilisation de LLMs (Large Language Models) pour rendre les dialogues du chatbot naturels.
    \item \textbf{Une infrastructure moderne :} Des bases de données vectorielles et une architecture Cloud pour assurer des temps de réponse instantanés.
\end{itemize}

\vspace{0.3cm}

% --- Synthèse ---
\section*{Synthèse}
La création de Wakala autour de ces deux composants (Chatbot et Moteur de Recommandation) garantit la création d'une expérience d'achat automobile innovante et sur-mesure. C'est l'exploitation des dernières avancées en Intelligence Artificielle générative au service direct du besoin utilisateur.

\end{document}
